\section{Traces and the absorption criterion}\label{sec:preliminaries}

For a unital \(\mathrm C^*\)-algebra \(D\), write \(\T(D)\) for its
tracial-state space. For \(\tau\in\T(D)\), set
\[
 \|x\|_{2,\tau}=\tau(x^*x)^{1/2}.
\]
Under the assumptions of Theorem~\ref{thm:main}, \(\T(A)\) is a
nonempty compact convex set. Indeed, nuclearity implies exactness,
and a unital stably finite exact algebra has a tracial state;
moreover its \(2\)-quasitraces are traces
\cite[Theorem~5.11 and Corollary~5.12]{Haagerup2014}.
Every tracial state on the simple algebra \(A\) is faithful:
its null space \(\{a:\tau(a^*a)=0\}\) is a closed two-sided ideal
that does not contain \(1\).

Fix a free ultrafilter \(\omega\) on \(\N=\{1,2,\ldots\}\). Define
\begin{equation}\label{eq:ultrapower}
 \|a\|_{2,u}=\sup_{\tau\in\T(A)}\|a\|_{2,\tau},\qquad
 A^\omega=\ell^\infty(\N,A)/I_\omega,\qquad
 I_\omega=\{(a_j):\lim_{j\to\omega}\|a_j\|_{2,u}=0\}.
\end{equation}
The inequalities
\(\|ab\|_{2,u}\leq\|a\|\|b\|_{2,u}\) and
\(\|ba\|_{2,u}\leq\|a\|\|b\|_{2,u}\) show that \(I_\omega\)
is a closed two-sided ideal. Faithfulness of traces embeds \(A\)
as the constant sequences in \(A^\omega\).
For any sequence \((\tau_j)\) in \(\T(A)\), the formula
\[
 \tau([(a_j)])=\lim_{j\to\omega}\tau_j(a_j)
\]
defines a tracial state on \(A^\omega\), since
\(\lvert\tau_j(a_j)\rvert\leq\|a_j\|_{2,\tau_j}\).
These are the \emph{limit traces}, denoted \(\T_\omega(A)\).
We do not identify them with all tracial states on \(A^\omega\).

Automorphisms permute \(\T(A)\) and therefore preserve
\(\|\cdot\|_{2,u}\). Applying \(\alpha_g\) to representative sequences
defines an action \(\bar\alpha:G\curvearrowright A^\omega\).
It preserves \(A^\omega\cap A'\). The quotient in
Equation~\eqref{eq:ultrapower} is also the quotient of the norm
ultrapower by its uniform trace-kernel ideal, since norm-null
sequences belong to \(I_\omega\).

\begin{definition}\label{def:gamma}
The action \(\alpha\) has \emph{equivariant uniform property Gamma}
if, for every integer \(k\geq2\), there are pairwise orthogonal
projections \(p_1,\ldots,p_k\in(A^\omega\cap A')^{\bar\alpha}\) with
\begin{equation}\label{eq:gamma}
 \tau(ap_i)=\frac{\tau(a)}{k}
 \qquad(a\in A,\ \tau\in\T_\omega(A),\ 1\leq i\leq k).
\end{equation}
\end{definition}

This is the unital simple formulation in
\cite[Definition~A]{SzaboWouters2025}.
Our construction also gives \(\sum_i p_i=1\).
We use the following established theorem.

\begin{theorem}[Szab\'o--Wouters]\label{thm:absorption}
Let a countable discrete amenable group act on a separable, unital,
simple, nuclear, \(\Zalg\)-stable \(\mathrm C^*\)-algebra with
nonempty tracial-state space. If the action has equivariant uniform
property Gamma, it is cocycle conjugate to its tensor product with
the trivial action on \(\Zalg\).
\end{theorem}

Theorem~\ref{thm:absorption} is the implication of
\cite[Corollary~C]{SzaboWouters2025} needed here.
Thus it remains to construct the projections required by Definition~\ref{def:gamma}.
The next two sections work inside a separable subalgebra \(C\) of
\(A^\omega\cap A'\). There we use \emph{all} traces in \(\T(C)\);
this will allow weighted restrictions of limit traces in
Section~\ref{sec:splittings}.
